\vfill\pagebreak
\section{Copies: \texttt{copy} and \texttt{dcopy}}
\label{sec:mutability:copies}

Sharing is the default, and a program that wants elements of its own asks for a
copy. \token{copy} was met as the way to make a slice from an array
(Section~\ref{sec:collections:copy}). It applies to a slice as well: it
reserves a new block, copies the elements into it, and gives a slice that
designates the new block. The next listing keeps the prices of
Listing~\ref{lst:mutability:discount} before the discount changes them.

\lstinputlisting[style=coloredverbatim, caption={Keeping the elements before a change, \textit{examples/mutability/backup.yr}}, label=lst:mutability:backup]{examples/mutability/backup.yr}
\vspace{-10pt}%
\begin{lstlisting}[style=bashVerb]
Now:    [108, 72, 40]
Seen:   [108, 72, 40]
Before: [120, 80, 45]
\end{lstlisting}

\token{seen}, on line~11, shares the elements of \token{prices}, and sees the
discount. \token{before}, on line~12, has elements of its own, and keeps the
prices as they were. A copy shares nothing with the original, so it needs no
\token{alias}, even to a variable declared \token{dmut}:
\token{let dmut other = copy prices;} is accepted, and changing \token{other}
leaves \token{prices} as it is.

A copy has a cost that sharing does not have: an allocation, and the time to
copy every element. A slice of a million elements passed to a function costs
two numbers, and its copy a million elements. A program copies when it needs
elements of its own, not by default.

\subsection{Copying a slice of slices}

The elements of a slice of slices are slices, and \token{copy} copies them as
slices are copied: a length and an address each. The copy has a new block of
rows, but its rows designate the elements of the rows of the original.
\token{dcopy}, for \textit{deep copy}, copies every level: the block of rows,
and each row with its elements.

\begin{lstlisting}[style=coloredverbatim]
let dmut grid = dcopy [[1, 2], [3, 4]];
let shallow = copy grid;
let dmut deep = dcopy grid;
grid[0][0] = 9;
deep[1][1] = 0;
println(grid, " ", shallow, " ", deep);
\end{lstlisting}
\vspace{-10pt}%
\begin{lstlisting}[style=bashVerb]
[[9, 2], [3, 4]] [[9, 2], [3, 4]] [[1, 2], [3, 0]]
\end{lstlisting}

\token{grid[0][0] = 9;} writes into the first row of \token{grid}, which is also
the first row of \token{shallow}. \token{deep} has rows of its own, and neither
sees the change of the other. Figure~\ref{fig:mutability:copies} shows the three
variables. The first line uses \token{dcopy} as well: \token{[[1, 2], [3, 4]]}
is an array of arrays, and \token{dcopy} turns each level into a slice, the
rows included.

\begin{figure}[h]
  \centering
  \begin{adjustbox}{max width=\linewidth}
    \begin{tikzpicture}
      \foreach \name/\y in {grid/0, shallow/-2.4, deep/-4.0} {
        \node[memcell] (\name len) at (0, \y) {2};
        \node[memaddr, right=0mm of \name len] (\name ptr) {};
        \node[memcell] (\name r0len) at (4.8, \y) {2};
        \node[memaddr, right=0mm of \name r0len] (\name r0ptr) {};
        \node[memcell, right=0mm of \name r0ptr] (\name r1len) {2};
        \node[memaddr, right=0mm of \name r1len] (\name r1ptr) {};
        \draw[mempoint] (\name ptr.center) -- (\name r0len.west);
      }
      \node[cflab, left=3mm of gridlen] (gridn) {\token{grid}};
      \node[cflab, left=3mm of shallowlen] (shallown) {\token{shallow}};
      \node[cflab, left=3mm of deeplen] (deepn) {\token{deep}};
      \node[memframe, fit=(gridn)(shallown)(gridptr)(deepptr)] (main) {};
      \node[memtitle] at ([xshift=2mm]main.north west) {main};

      \node[memcell] (a0) at (5.8, -1.2) {1};
      \node[memcell, right=0mm of a0] (a1) {2};
      \node[memcell] (b0) at (8.3, -1.2) {3};
      \node[memcell, right=0mm of b0] (b1) {4};
      \node[memcell] (c0) at (5.8, -5.2) {1};
      \node[memcell, right=0mm of c0] (c1) {2};
      \node[memcell] (e0) at (8.3, -5.2) {3};
      \node[memcell, right=0mm of e0] (e1) {4};

      \draw[mempoint] (gridr0ptr.center) to[out=-90, in=90] (a0.north);
      \draw[mempoint] (gridr1ptr.center) to[out=-90, in=90] (b0.north);
      \draw[mempoint] (shallowr0ptr.center) to[out=90, in=-90] (a0.south);
      \draw[mempoint] (shallowr1ptr.center) to[out=90, in=-90] (b0.south);
      \draw[mempoint] (deepr0ptr.center) to[out=-90, in=90] (c0.north);
      \draw[mempoint] (deepr1ptr.center) to[out=-90, in=90] (e0.north);

      \node[memregion] (stack) at ($(main.north) + (0, 5mm)$) {stack};
      \node[memregion] at ({$(gridr0len)!0.5!(b1)$} |- stack.base) {heap};
      \draw[memsep] (3.3, 1.2) -- (3.3, -5.8);
    \end{tikzpicture}
  \end{adjustbox}
  \caption{\label{fig:mutability:copies} A slice of two rows, \token{grid}, and
    two copies of it. \token{shallow = copy grid} has a block of its own for
    the two rows, but each row is copied as a slice is, length and address,
    and still designates the elements of \token{grid}. \token{deep = dcopy grid}
    copies the rows as well: it shares nothing with \token{grid}.}
\end{figure}


Since the rows of \token{shallow} are those of \token{grid}, whose elements can
change, \token{shallow} cannot be declared \token{dmut}. The type of
\token{copy grid} is \token{mut [mut [i32]]}: the rows are constant through the
copy, and \token{let dmut shallow = copy grid;} is refused with
\errmsg{E4045}{discarding the constant property is prohibited}, the note under
it saying \textit{from mut [mut [i32]] to mut [mut [mut i32]]}. A copy made
with \token{copy} can therefore never change the rows of the original. To
change the rows of a copy, the copy is made with \token{dcopy}.

\subsection{How deep \texttt{dcopy} goes}

\token{dcopy} is defined in terms of itself. The deep copy of a slice is a new
block whose elements are the deep copies of the elements of the slice. When the
elements are slices in turn, each is deep-copied the same way, and so on, until
the elements borrow nothing: an integer is its own deep copy. Written by hand
for a grid, one level of that recursion is a \token{copy} of each row.

\begin{lstlisting}[style=coloredverbatim]
fn deepCopy(grid: [[i32]])-> dmut [[i32]] {
  let dmut result: [[i32]] = [];
  for row in grid {
    result ~= [copy row];
  }
  alias result
}
\end{lstlisting}

\token{dcopy} does the same at every level, however many there are. A cube of
numbers, a slice of grids of type \token{[[[i32]]]}, has three levels, and
\token{dcopy cube} copies the block of grids, the block of rows of each grid,
and the elements of each row.

\begin{lstlisting}[style=coloredverbatim]
let dmut cube = dcopy [[[1, 2], [3, 4]], [[5, 6], [7, 8]]];
let dmut other = dcopy cube;
other[1][1][1] = 0;
println(cube, " ", other);
\end{lstlisting}
\vspace{-10pt}%
\begin{lstlisting}[style=bashVerb]
[[[1, 2], [3, 4]], [[5, 6], [7, 8]]] [[[1, 2], [3, 4]], [[5, 6], [7, 0]]]
\end{lstlisting}

\subsubsection*{A value that contains itself}

A recursion needs a base case (Section~\ref{sec:functions:infinite_recursion}),
and that of \token{dcopy} is the element that borrows nothing. A slice that
contained itself, directly or through another slice, would never reach it: the
deep copy of its elements would include its own deep copy, which would include
it again, and so on until the stack overflows. The types prevent such a slice
from existing.

\begin{lstlisting}[style=coloredverbatimError, escapechar=@, caption=A slice that would contain itself]
use std::io;

fn main() {
  let dmut b = copy [copy [1]];
  let dmut a = copy [alias b];
  b[0] = @\hb{a}@; // error
  println(a);
}
\end{lstlisting}

The compiler says \errmsg{E4095}{incompatible types mut [mut i32] and mut [mut
  [mut [mut i32]]]}. \token{b} is a slice of slices of integers, and \token{a}
a slice of such slices: one level more. For \token{b[0] = a;} to be accepted,
an element of \token{b} would have to have the type of \token{a}, which holds
the type of \token{b}, which holds that of its elements, without end. The type
of a slice is always one level deeper than that of its elements, so no slice,
array or map can hold itself, and \token{dcopy} always reaches the bottom.

The classes of Chapter~\ref{chap:custom_types} can refer to each other, and to
themselves through others: two objects each naming the other form a cycle. The
compiler therefore does not deep-copy a class by itself: \token{dcopy} on an
object whose class does not say how it is copied is refused with
\errmsg{E4018}{class type \errhole{} does not implement trait
  core::types::Copiable}, and the class decides what its deep copy does with a
cycle.

\subsubsection*{Maps and tuples}

A map is deep-copied the same way, its values included: the deep copy of a map
from names to slices of grades has slices of its own. Tuples are the exception.
Neither \token{copy} nor \token{dcopy} applies to a tuple, which is refused with
\errmsg{E4167}{no copy exists for type (i32, [i32])}, and therefore not to a
slice whose elements are tuples holding slices either. Such a value is copied
field by field.

\subsection{Who makes the copy}

A function that changes a slice can be written in two ways. The first changes
the elements it receives, in place, and the second copies them first and
returns the copy, changed. The next listing sorts a slice in increasing order,
in place: it looks for the smallest element and exchanges it with the first,
then for the smallest of the others, exchanged with the second, and so on. This
method is called the \textit{selection sort}.

\lstinputlisting[style=coloredverbatim, caption={Sorting in place, \textit{examples/mutability/sort.yr}}, label=lst:mutability:sort]{examples/mutability/sort.yr}
\vspace{-10pt}%
\begin{lstlisting}[style=bashVerb]
[3, 7, 19, 25, 42]
[42, 7, 19, 3, 25] [3, 7, 19, 25, 42]
\end{lstlisting}

\token{sort} takes the slice as a \token{dmut} parameter, and makes no copy.
Its caller chooses. On line~23, \token{main} sorts \token{scores} itself, and
pays for no copy. On lines~27 and~28, it wants to keep \token{arrivals} in the
order of arrival, so it makes the copy, and sorts the copy.

The other way puts the copy inside the function.

%% check: skip
\begin{lstlisting}[style=coloredverbatim]
fn sorted(values: [i32])-> [i32] {
  let dmut result = copy values;
  sort(alias result);
  result
}
\end{lstlisting}

\token{sorted} never changes its argument, but every call copies it, including
\token{scores = sorted(scores);}, which throws the old elements away at once.
The caller of \token{sorted} cannot avoid the copy, whereas the caller of
\token{sort} makes one only when it needs one. A function that changes
elements is therefore best written in place, with a \token{dmut} parameter, and
the decision to copy left to the caller, who knows whether the original is
still needed.

The copy belongs inside the function when its contract says that it only reads:
a function that computes the median of a slice, the element in the middle once
the slice is sorted, takes a plain parameter, and must not sort the slice of its
caller. It sorts a copy, as Exercise~8 does.

\subsection{The repeated row}
\label{sec:mutability:repeated_row}

A caution of Section~\ref{sec:collections:grids} warned against
\token{copy [copy [0 ; cols] ; rows]}, which looks like a grid of \token{rows}
rows but changes all of them at once. The reason is now visible. \token{[v ; n]}
computes \token{v} once, and makes \token{n} copies of the value: for a slice,
\token{n} copies of the same length and the same address.
Figure~\ref{fig:mutability:repeated_row} shows the grid.

\begin{figure}[h]
  \centering
  \begin{adjustbox}{max width=\linewidth}
    \begin{tikzpicture}
      \node[memcell] (glen) at (0, 0) {3};
      \node[memaddr, right=0mm of glen] (gptr) {};
      \node[cflab, left=3mm of glen] (gn) {\token{grid}};
      \node[memframe, fit=(gn)(gptr)] (main) {};
      \node[memtitle] at ([xshift=2mm]main.north west) {main};

      \node[memcell] (r0len) at (4.8, 0) {4};
      \node[memaddr, right=0mm of r0len] (r0ptr) {};
      \node[memcell, right=0mm of r0ptr] (r1len) {4};
      \node[memaddr, right=0mm of r1len] (r1ptr) {};
      \node[memcell, right=0mm of r1ptr] (r2len) {4};
      \node[memaddr, right=0mm of r2len] (r2ptr) {};
      \draw[mempoint] (gptr.center) -- (r0len.west);
      \foreach \i in {0, 1, 2} {
        \node[memnote, above=0.5mm of r\i len.north east] {row \i};
      }

      \foreach \value [count=\i from 0] in {0, 0, 7, 0} {
        \ifnum\i=2
          \node[memhit] (e\i) at (7.2 + \i * 1.0, -1.9) {\value};
        \else
          \node[memcell] (e\i) at (7.2 + \i * 1.0, -1.9) {\value};
        \fi
      }
      \foreach \i in {0, 1, 2} {
        \draw[mempoint] (r\i ptr.center) to[out=-90, in=90] (e0.north);
      }

      \node[memregion] (stack) at ($(main.north) + (0, 5mm)$) {stack};
      \node[memregion] at ({$(r0len)!0.5!(r2ptr)$} |- stack.base) {heap};
      \draw[memsep] (3.3, 0.9) -- (3.3, -2.6);
    \end{tikzpicture}
  \end{adjustbox}
  \caption{\label{fig:mutability:repeated_row} The grid
    \token{copy [copy [0 ; 4] ; 3]}. \token{copy [0 ; 4]} made one row, and
    \token{[v ; 3]} repeated that value three times: three copies of the same
    length and the same address. \token{grid[1][2] = 7;} changes the element
    that the three rows designate.}
\end{figure}


\begin{lstlisting}[style=coloredverbatim]
let dmut grid = copy [copy [0 ; 4] ; 3];
grid[1][2] = 7;
println(grid);
\end{lstlisting}
\vspace{-10pt}%
\begin{lstlisting}[style=bashVerb]
[[0, 0, 7, 0], [0, 0, 7, 0], [0, 0, 7, 0]]
\end{lstlisting}

The outer \token{copy} copies the three rows as slices, and does not separate
them. The loop of Section~\ref{sec:collections:grids} makes a new row at each
turn, with a \token{copy} of its own, and gives distinct rows. \token{dcopy} of
the grid above also gives distinct rows, since it copies each of them in turn:
\token{let dmut fixed = dcopy grid;} is a grid of three rows of its own, which
start with the 7 of the shared row. In version \gycversion{} of the compiler,
\token{dcopy} written directly around the repeated literal,
\token{dcopy [copy [0 ; 4] ; 3]}, still gives a single shared row, and the loop
is the way to write it.
